\appendix
\section{A published-input density lemma for polylogarithmic conductors}
\label{app:pzd}

The all-center argument only uses conductors bounded by a fixed power of the
logarithm. The following narrower lemma avoids any appeal to an unrefereed
hybrid character-aspect theorem while retaining the $30/13$ coefficient.

For a Dirichlet character $\chi$, let $N(\sigma,T,\chi)$ count, with
multiplicity, the zeros $\rho=\beta+i\gamma$ of $L(s,\chi)$ satisfying
$\beta\ge\sigma$ and $|\gamma|\le T$.

\begin{theorem}[Polylogarithmic-conductor density]\label{thm:pzd}
Fix $K>0$, $\delta>0$, and $\eta>0$.  Uniformly for
\[
 T\ge T_0(K,\delta,\eta),\qquad
 \frac12+\delta\le\sigma\le\frac45,
 \qquad Q\le(\log T)^K,
\]
one has
\[
 \sum_{q\le Q}\sum_{\chi\bmod q}N(\sigma,T,\chi)
 \ll_{K,\delta,\eta}
 T^{\frac{30}{13}(1-\sigma)+\eta}.
 \tag{A.1}
\]
\end{theorem}

\begin{proof}
We first reduce to primitive characters.  If $\chi\pmod q$ is induced by the
primitive character $\chi^*\pmod r$, then
\[
 L(s,\chi)=L(s,\chi^*)
 \prod_{\substack{p\mid q\\p\nmid r}}
       (1-\chi^*(p)p^{-s}).
\]
Every zero of a factor in the finite product has real part zero.  Consequently
the zeros counted in (A.1) are zeros of $L(s,\chi^*)$.  The number of
character--modulus pairs under consideration is $O(Q^2)$, so it is enough to
prove, for every primitive $\chi\pmod r$, $r\le Q$,
\[
 N(\sigma,T,\chi)
 \ll_{K,\delta,\eta}
 T^{\frac{30}{13}(1-\sigma)+\eta/2}.
 \tag{A.2}
\]

We give the reduction of (A.2) to published large-values and moment estimates.
This also records where conductor uniformity enters. Fix small constants
$\kappa>0$ and $0<\vartheta<1$, to be chosen at the end in terms of $\eta$, and put
\[
 \mathcal R=r(T+2),\qquad U=2\mathcal R^\kappa,
 \qquad Y=\mathcal R^{1/2}.
\]
If $r=1$, the character is the zeta character.  The range
$7/10\le\sigma\le4/5$ follows from \cite[Theorem~1.2]{GuthMaynard2026},
while the lower range follows from the Ingham estimate summarized in
\cite[(1.4)]{GuthMaynard2026}.  We may therefore assume $r>1$;
the primitive character $\chi$ is then nonprincipal and no pole residue
occurs. Form the usual mollifier
\[
 M_U(s,\chi)=\sum_{d\le U}\mu(d)\chi(d)d^{-s}.
\]
For $\operatorname{Re} s>1$,
\[
 L(s,\chi)M_U(s,\chi)
 =\sum_{n\ge1}c_n\chi(n)n^{-s},
 \qquad c_n=\sum_{\substack{d\mid n\\d\le U}}\mu(d).
 \tag{A.3}
\]
Thus $c_1=1$ and $c_n=0$ for $1<n\le U$. We now record the detector without
discarding its weights. Mellin inversion for $e^{-n/Y}$ and a contour shift
to $\operatorname{Re} z=1/2-\beta$ give, at a zero
$\rho=\beta+i\gamma$,
\begin{align}
 e^{-1/Y}+\sum_{n>U}c_n\chi(n)n^{-\rho}e^{-n/Y}
  ={}&\frac1{2\pi i}\int_{\operatorname{Re} z=1/2-\beta}
       \Gamma(z)Y^zL(\rho+z,\chi)M_U(\rho+z,\chi)\,dz.         \tag{A.4}
\end{align}
The apparent pole at $z=0$ is removable: since $L(\rho,\chi)=0$,
$L(\rho+z,\chi)=zH_\rho(z)$ with $H_\rho$ holomorphic near zero, while
$z\Gamma(z)=\Gamma(z+1)$.  There is no pole at $z=1-\rho$ because $\chi$
is nonprincipal.  If
$N_\chi(t,t+1)$ denotes the number of nontrivial zeros in the critical strip,
counted with multiplicity, whose ordinates lie in $[t,t+1]$, we use the local
zero count
\[
 N_\chi(t,t+1)\ll\log(r(|t|+2)).                                \tag{A.5}
\]
This estimate controls multiplicity.  Put
$a=1/2-\beta$ and $B=(\log\mathcal R)^2$.  Parametrizing the shifted line by
$z=a+iv$ gives $dz=i\,dv$ and $\rho+z=1/2+i(\gamma+v)$.  Equivalently, with
the absolute critical-line ordinate $u=\gamma+v$, the integrand is
\[
 \Gamma\bigl(a+i(u-\gamma)\bigr)
 Y^{a+i(u-\gamma)}L(1/2+iu,\chi)M_U(1/2+iu,\chi),
\]
and the truncated interval is $|u-\gamma|\le B$.  Uniformly for
$1/2<\beta\le1$ and $|v|\ge1$, the fixed-strip bound for $L$, the elementary
bound $|M_U(1/2+i(\gamma+v),\chi)|\le U+1$, and vertical Gamma decay give
\begin{align*}
 &\left|\Gamma(a+iv)Y^{a+iv}
 L(1/2+i(\gamma+v),\chi)M_U(1/2+i(\gamma+v),\chi)\right| \\
 &\qquad\ll r^2(U+1)(1+|v|)(5+|\gamma|+|v|)^2e^{-|v|} \\
 &\qquad\ll r^2(U+1)(5+|\gamma|)^2(1+|v|)^3e^{-|v|}.
\end{align*}
Hence the two vertical tails are
\[
 \ll r^2(U+1)(5+|\gamma|)^2e^{-B/2}.
\]
If $N=\lceil YB\rceil$, the arithmetic tail is at most
\[
 (U+1)e^{-(N+1)/Y}(1-e^{-1/Y})^{-1}
 \ll (U+1)Y e^{-B}.
\]
Since $\mathcal R=r(T+2)$, $Y=\mathcal R^{1/2}$,
$U=2\mathcal R^\kappa$, $|\gamma|\le T$, and
$r\le(\log T)^K$, both tails are $O_A(\mathcal R^{-A})$ for every fixed
$A$. Consequently every remaining zero satisfies either
\[
 \left|\sum_{U<n\le Y(\log\mathcal R)^2}
 c_n\chi(n)n^{-\rho}e^{-n/Y}\right|\gg1,                        \tag{I}
\]
or the shifted critical-line integral containing
$L(1/2+i(\gamma+u),\chi)M_U(1/2+i(\gamma+u),\chi)$ is large for
$|u|\le(\log\mathcal R)^2$. This is the classical Montgomery
Dirichlet-$L$ detector \cite[Chapters~10 and 12]{Montgomery1971}; (A.3) displays
the character-twisted coefficient identity. Equation (A.5) controls
multiplicity and permits thinning at any prescribed separation, with the
separation length times a logarithm as the loss.

For the critical-line alternative, first thin the original ordinates at
separation $3(\log\mathcal R)^2$. The shifted ordinates are then
one-separated, and the thinning loss is polylogarithmic by (A.5). The detector
inequality and the integrability of the gamma factor imply, for one shifted
ordinate $v$, that
\[
 |L(1/2+iv,\chi)M_U(1/2+iv,\chi)|
 \gg Y^{\sigma-1/2}\mathcal R^{-o(1)}.
\]
Since
\[
 |M_U(1/2+iv,\chi)|\le\sum_{d\le U}d^{-1/2}\ll U^{1/2},
\]
Montgomery's discrete family fourth moment
\cite[Theorem~10.3]{Montgomery1971}, restricted to the fixed primitive character
by nonnegativity,
\[
 \sum_v|L(1/2+iv,\chi)|^4\ll(rT)^{1+o(1)}
\]
at one-separated ordinates gives
\[
 \#\mathcal Z_{II}\ll\mathcal R^{2(1-\sigma)+2\kappa+o(1)}.           \tag{A.6}
\]
All conductor factors are harmless because
$r\le(\log T)^K$. Notice that we used the pointwise bound for $M_U$; a
fourth moment of $L$ and only a second moment of $M_U$ would not justify (A.6).

It remains to count the zeros satisfying (I). Dyadic subdivision supplies
one $N$ with
\[
 \mathcal R^\kappa\ll N\ll\mathcal R^{1/2}
                    (\log\mathcal R)^2
 \tag{A.7}
\]
and a large block for a proportion $(\log\mathcal R)^{-O(1)}$ of these
zeros. The lower bound holds because the coefficients in (A.3) vanish for
$1<n\le U$.
We retain the exponential weight in the coefficients.  The real part of a
zero must also be removed before applying one large-values theorem to the
whole set. Fix the left edge $\sigma$, put $L=\log N$, choose a fixed
$\zeta_0\in C_c^\infty((1/2,2))$ equal to one near
$[1,1+\log(2)/L]$, and set
\[
 \psi_\beta(u)=\zeta_0(u/L)e^{-(\beta-\sigma)u}.
\]
Uniformly for $\sigma\le\beta\le1$,
\[
 \|\psi_\beta^{(j)}\|_1\ll_jL^{1-j},\qquad
 |\widehat\psi_\beta(\xi)|\ll_jL^{1-j}|\xi|^{-j},
 \qquad \|\widehat\psi_\beta\|_1\ll1.                         \tag{A.8}
\]
Fourier inversion gives, up to the fixed Fourier-sign convention,
\[
 \sum_{N<n\le2N}c_ne^{-n/Y}\chi(n)n^{-\beta-i\gamma}
 =\int_{\mathbb R}\widehat\psi_\beta(\xi)
   \sum_{N<n\le2N}c_ne^{-n/Y}\chi(n)
        n^{-\sigma-i(\gamma+2\pi\xi)}\,d\xi.
 \tag{A.9}
\]
Put $H=\mathcal R^\vartheta$. By (A.8), choosing the
integration-by-parts order in terms of $\vartheta$ makes the contribution of
$|\xi|>H$ smaller than the block threshold, uniformly in $\beta$. Indeed,
$|c_n|\le d(n)$ gives the crude uniform estimate
\[
 \sup_t\left|
 \sum_{N<n\le2N}c_ne^{-n/Y}\chi(n)n^{-\sigma-it}
 \right|
 \ll \mathcal R^{o(1)}N^{1-\sigma},
\]
and its product with the Fourier tail from (A.8) is smaller than any prescribed
fixed power of $\mathcal R^{-1}$ after increasing the integration-by-parts
order. Thus every
retained zero chooses some $|\xi_\rho|\le H$ where a common polynomial is
large. Normalize with
\[
 C_N=\max_{N<n\le2N}|c_n|=\mathcal R^{o(1)}
\]
(the block is nonzero), and define the zero-independent polynomial
\[
 D_\chi(t)=\sum_{N<n\le2N}a_n\chi(n)n^{it},
 \qquad a_n=\frac{c_ne^{-n/Y}(N/n)^\sigma}{C_N},
 \qquad |a_n|\le1,
 \tag{A.10}
\]
which satisfies
\[
 |D_\chi(t_j)|\ge N^\sigma\mathcal R^{-o(1)}
 \tag{A.11}
\]
at a shifted ordinate $t_j=-\gamma-2\pi\xi_\rho$. Any unit interval contains
at most $O((H+1)\log\mathcal R)$ of these shifted ordinates, counted with
multiplicity: all their original ordinates lie in an interval of length
$O(H)$, so this follows from (A.5). Greedy selection therefore leaves a
one-separated subset of cardinality
\[
 \gg\#\mathcal Z_I/((H+1)\log\mathcal R).                             \tag{A.12}
\]
It lies in an interval of length $O(T+H)=O(T)$. Translating that interval
multiplies the common coefficients by unimodular factors. This quantifies the
real-part-removal step compressed in \cite[Section~13.1]{GuthMaynard2026}; the fixed factor
$\chi(n)e^{-n/Y}$ remains in the coefficients throughout.

Guth and Maynard's published large-values theorem \cite[Theorem~1.1]{GuthMaynard2026}
applies to every complex coefficient sequence bounded by one.  Hence it
applies to (A.10).  Their power argument in
\cite[Section~13.1]{GuthMaynard2026} has two branches: Theorem~1.1 is used
when the powered length is at most their switching scale, and the usual
discrete mean-value theorem is used in the complementary branch.  Both inputs
allow arbitrary bounded complex coefficients.  In the present twisted setting,
complete multiplicativity gives
\[
 \chi(n_1)\cdots\chi(n_k)=\chi(n_1\cdots n_k),
\]
and, literally,
\[
 D_\chi^k(t)=
 \sum_{N^k<m\le(2N)^k}\chi(m)b_m m^{it},
 \qquad |b_m|\le d_k(m)=\mathcal R^{o(1)}.
\]
Split this support into $O_k(1)$ dyadic blocks. The triangle inequality and
pigeonholing retain one common block for a fixed positive proportion of the
selected ordinates; normalize its divisor-bounded coefficients before
applying \cite[Theorem~1.1]{GuthMaynard2026}. These operations cost only
$\mathcal R^{o(1)}$. By (A.7) and $r\le(\log T)^K$, the polynomial length obeys
$N\le T^{1/2+o_K(1)}$ and $N\ge T^{\kappa+o_K(1)}$. The integer $k$ in
their proof is therefore bounded in terms of $\kappa$; this is enough because
$\kappa$ is fixed before $T\to\infty$. The same power argument gives, for
$7/10\le\sigma\le4/5$,
\[
 \#\mathcal Z_I
 \le T^{\frac{15(1-\sigma)}{3+5\sigma}
          +\vartheta+o_{K,\kappa}(1)}.
 \tag{A.13}
\]
Because
\[
 \frac{15}{3+5\sigma}\le\frac{30}{13}
 \quad\left(\sigma\ge\frac7{10}\right),
 \qquad
 \min_{7/10\le\sigma\le4/5}
 \left(\frac{15}{3+5\sigma}-2\right)(1-\sigma)
 =\frac1{35},
\]
we may first choose $\kappa$ so that $2\kappa<1/70$ and then choose
$\vartheta$ and all remaining subpower losses smaller than the allocated
$\eta/2$. Equations (A.6) and (A.13) prove (A.2) in this range.

For $1/2+\delta\le\sigma\le7/10$, Montgomery's classical family estimate
\cite[Theorem~12.1]{Montgomery1971} gives, for each modulus $r$,
\[
 \sum_{\chi\bmod r}N(\sigma,T,\chi)
 \ll (rT)^{\frac{3(1-\sigma)}{2-\sigma}}
       (\log(rT))^9;
 \tag{A.14}
\]
for the fixed character under consideration, nonnegativity gives explicitly
\[
 N(\sigma,T,\chi)
 \le \sum_{\psi\bmod r}N(\sigma,T,\psi).
\]
Since
\[
 \frac3{2-\sigma}\le\frac{30}{13}
 \quad\left(\sigma\le\frac7{10}\right),
\]
(A.2) follows throughout the asserted strip after taking $T$ sufficiently
large in terms of $K,\delta,\eta$.

Finally $Q^{O(1)}=(\log T)^{O_K(1)}=T^{o_K(1)}$.  Summing (A.2) over the
$O(Q^2)$ original character--modulus pairs and allocating the total subpower
loss to $\eta/2$ proves (A.1).
\end{proof}

\begin{remark}[Scope of the density lemma]
The theorem does not extend the Guth--Maynard estimate to a character family.
It applies that published arbitrary-coefficient theorem separately to each
character and uses that the family is only polylogarithmic. The fixed-character
detector, fourth-moment input, real-part removal, spacing loss, and bounded-power
coefficient normalization are displayed above because they are the uniformity
points on which this reduction depends.
\end{remark}
